\section{Definitions and Threat Model}
\label{sec:definitions}

\subsection{Purpose and authorization}

Let a \emph{case} be a tuple of observable attributes $x$ (public features) and
confidential attributes $z$. A \emph{purpose} $p$ is a task label under which a
system may act on a case (e.g., \textsc{hardship-support-routing},
\textsc{fraud-review}). For a confidential attribute $z_i$, an authorization
mapping $\mathrm{perm}(z_i, p) \in \{\textsc{auth},\textsc{proh}\}$ states
whether the organization may use $z_i$ when acting under purpose $p$. A
\emph{decision function} $d(\tilde x, p) \to \mathcal{A}$ maps the rendered
input $\tilde x$ (a prompt over visible attributes) and the purpose to an
action in $\mathcal{A}$.

\subsection{The PSBE property}

For a fixed case and fixed public features, consider three renderings of the
same confidential attribute $z_i$: visible ($\tilde x^{+}$), masked
($\tilde x^{-}$), and the identity of the model's own nondeterminism (two
identical renderings). We define

\begin{itemize}
  \item \textbf{Authorized utility} $U(p)$: decision quality on the authorized
  task, e.g., balanced accuracy (BACC) against a ground-truth label.
  \item \textbf{Influence} $I(p)$: the rate at which a decision \emph{changes}
  when a masked attribute becomes visible, under purpose $p$:
  $I(p) = \Pr\bigl[d(\tilde x^{+}, p) \ne d(\tilde x^{-}, p)\bigr]$.
  \item \textbf{Natural-decision floor} $F$: the rate at which the decision
  changes when the \emph{identical} input is rendered twice,
  $F = \Pr\bigl[d(\tilde x, p) \ne d'(\tilde x, p)\bigr]$ over repeated calls.
  Nondeterministic LLM sampling makes $F>0$.
\end{itemize}

\textbf{Definition (PSBE).} A system is \emph{purpose-selective} on attribute
$z_i$ with margin $\delta$ and retention gate $\rho$ if, for every authorized
purpose $p^{+}$ and prohibited purpose $p^{-}$:
\begin{equation}
\begin{aligned}
I_{\mathrm{governed}}(p^{-}) &\le F + \delta &&\text{forbidden to influence}\\
U_{\mathrm{governed}}(p^{+}) &\ge \rho \cdot U_{\mathrm{ungoverned}}(p^{+}) &&\text{authorized to use}.
\end{aligned}
\label{eq:psbe}
\end{equation}
The first conjunct bounds prohibited influence at the floor (the system is
\emph{forbidden to influence}); the second requires the authorized utility to
survive governance (the system remains \emph{authorized to use}). Both are
behavioral, decision-level properties: they hold or fail at the model output,
regardless of where data was stored or transmitted.

\subsection{Threat model}

We target \emph{operational purpose confusion}---an organization that holds a
field for a lawful purpose and feeds it to a model that acts under another
purpose---not adversarial leakage or prompt injection. Concretely: the threat
model includes negligent or erroneous deployment (field included when it should
not be; no runtime guard), and honest-but-curious model providers (the operator
controls the model call through its own governed adapter). It explicitly
excludes: (i) prompt-injection attacks that attempt to exfiltrate the field;
(ii) a hostile provider that colludes to train on the field (no training-data
protection is claimed); (iii) formal non-interference guarantees (we measure,
we do not prove). We additionally require that the deployed system never
automatically falls back to an ungoverned path when the governed path fails,
because a fallback would silently bypass the property in Equation~\eqref{eq:psbe}.
